\documentclass[b5paper,11pt]{article}
\usepackage{notestyle}
\usepackage{amsmath}
\usepackage{paracol}

% 避免在未显示编号时出现 subsection 书签锚点重复
\renewcommand{\theHsection}{section.\arabic{section}}
\renewcommand{\theHsubsection}{\theHsection.\arabic{subsection}}
\renewcommand{\theHsubsubsection}{\theHsubsection.\arabic{subsubsection}}
\renewcommand{\theHparagraph}{\theHsubsubsection.\arabic{paragraph}}

% 文档专属：页眉文字
\fancyhead[L]{\fontsize{8pt}{10pt}\selectfont\color{gray!70}教育学通论}

% 添加右下角页脚页数标识
\fancyfoot[R]{\fontsize{8pt}{10pt}\selectfont\color{gray!70} \thepage}

% 文档信息
\title{教育学通论}
\date{2026/04/12}
\author{ClaudeRainer}

\begin{document}
\maketitle

\section{第一章 教育的发展}
\subsection{第一节 \space 教育概述}
\subsubsection{教育的概念}
“教育”最早出现在《孟子·尽心上》“得天下英才而教育之，三乐也”一句中，是指培养人才的过程。现代教育学家对教育的定义有多种表述，但核心内容基本一致：教育是指在一定的社会条件下，通过教师和学生之间的互动，传递知识、技能、价值观和文化传统的过程。

东汉许慎在《说文解字》中说：“教，上所施，下所效也。”这表明教育是一种培养人的活动，目的在于使人为善。

\paragraph{教育的定义}
教育分为\code{广义教育}和\code{狭义教育}两种定义：
\begin{itemize}
  \item 广义教育包括：社会教育、学校教育、家庭教育等。
  \item 狭义教育特指学校教育，\textcolor{ctp-red}{有目的、有计划、有组织}地对受教育者的\code{身心}施加影响。
\end{itemize}

狭义教育是把受教育者培养成\textcolor{ctp-red}{一定社会所需要的人}的活动。因此狭义教育具有如下几个基本内涵：
\begin{itemize}
  \item \code{受社会制约}：教育必须适应社会发展的需要，培养符合社会需求的人才。
  \item \code{有目的、有计划、有组织}：教育活动必须有明确的目标，科学的计划和有效的组织。
  \item \code{双边活动}：教育是由教育者和受教育者共同参与的双边活动。
\end{itemize}

\subsubsection{教育的构成要素}
教育要素是指在教育活动和教育过程中必不可少的、最基本的构成因素，主要由以下几个要素构成：
\textcolor{ctp-peach}{教育者、受教育者、教育内容、教育场所}

\paragraph{教育者}
在教育活动中指导、影响受教育者学习与素质发展的人。具有如下特点：
\begin{itemize}
  \item \code{职业性}：具有一定的权威性，是知识传播和教育活动的承担者。
  \item \code{专业性}：具有专业的训练、专业工作的资格、专业素质、享有一定社会地位。
\end{itemize}
基于以上这些特点，教育者在教育活动中应起到\textcolor{ctp-red}{引导、组织、指导}等主导作用。

\paragraph{受教育者}
指在教育活动中，接受\textcolor{ctp-red}{教育指导和影响}的人，具有如下特点：
\begin{itemize}
  \item \textcolor{ctp-red}{发展性和不成熟性}： 受教育者处于不断发展变化的过程中，具有成长的潜力和不成熟的特点。
  \item \textcolor{ctp-red}{可塑性和可教性}：\textcolor{ctp-red}{只要教育得当，每个人都可以得到发展。}
  \item \textcolor{ctp-red}{能动性和主体性}：\textcolor{ctp-rosewater}{有一定的知识经验和思想意识}，并具有\textcolor{ctp-rosewater}{一定的动机。}
\end{itemize}

\vspace{0.6em}

因此，受教育者既是\textcolor{ctp-red}{教育的对象}，也是\textcolor{ctp-red}{学习活动的主体}。
相对于教育者来说，受教育者处于被引导的地位；
相对于学习对象来说，\textcolor{ctp-red}{受教育者是学习的主体}。

\paragraph{教育内容}
\textcolor{ctp-red}{教育内容是教育者和受教育者共同加工的对象。}

\begin{itemize}
  \item \textcolor{ctp-red}{系统性和价值性} 蕴含一定社会或统治阶级的思想以及一定价值观。
  \item \textcolor{ctp-red}{发展性与开放性} 知识是人类社会的智慧结晶和经验总结，是不断发展的。
  \item \textcolor{ctp-red}{广泛性与载体的多样性} 教育内容的范围非常广泛，载体也多种多样。
\end{itemize}

\paragraph{教育场所}
教育活动发生的具体地点，包括学校、家庭、社会等，是教育过程的重要组成部分。
教育场所具有\textcolor{ctp-red}{多样性和复杂性}，不同的教育场所对教育活动的组织和实施有着不同的影响。

\subsubsection{教育的发展过程}
教育的发展大体经历了\textcolor{ctp-red}{非形式化教育、形式化教育、制度化教育}三个阶段。
\begin{itemize}
  \item \textcolor{ctp-red}{非形式化教育}：指没有\textcolor{ctp-maroon}{固定的教育者和受教育者}，与生活过程、生产过程混为一体的教育方式。（一直存在，直到原始社会的解体）
  \item \textcolor{ctp-red}{形式化教育}：有\textcolor{ctp-maroon}{相对固定的教育者和受教育者}，但没有形成制度化的教育体系。\textcolor{ctp-teal}{从生产环境中独立出来学习。}（从原始社会末期开始）
  \item \textcolor{ctp-red}{制度化教育}：形成了相对完善的教育体系和制度，教育活动有明确的规范和组织形式。（从奴隶社会开始）
\end{itemize}

\subsection{第二节 \space 中国教育的发展}

\subsubsection{中国古代教育}
\textcolor{ctp-red}{我国是历史上最早}出现学校教育的国家，大约在公元前\textcolor{ctp-red}{3000年}的原始社会末期出现了学校的萌芽。
夏、商、周时期\textcolor{ctp-red}{形式化教育}开始，后经汉、唐、宋等朝代发展，进一步发展和完善。
\paragraph{学校教育制度} 中国古代学校教育按其性质分为\textcolor{ctp-red}{官学、私学、书院}三种形式。

\subparagraph{官学}
凡有\textcolor{ctp-red}{官府举办}管辖的学校称为官学，由朝廷直接办理的学校为\textcolor{ctp-red}{中央官学}，
如西周国学，汉代太学，唐朝国子学，元、明、清的国子监等；由地方官府举办的\textcolor{ctp-red}{学校}，如西周乡学、汉代郡国学、唐代府州县学、
元代路学及社学、明清府州县学及卫学。

\textcolor{ctp-red}{西周}建立了\textcolor{ctp-red}{比较完备}的学校教育制度。学校分为\textcolor{ctp-red}{国学和乡学}两种。
国学专为\textcolor{ctp-red}{奴隶主贵族子弟}设立，按入学年龄和程度分为\textcolor{ctp-red}{小学，大学}两级；
乡学则为\textcolor{ctp-red}{奴隶主子弟和庶民中优秀者}设立，只有\textcolor{ctp-red}{小学}一级。

\textcolor{ctp-red}{汉唐}逐渐建立\textcolor{ctp-red}{中央官学和地方官学}的学校教育体制。公元前141年，\textcolor{ctp-red}{文翁兴学}。

\textcolor{ctp-red}{唐朝}建立了\textcolor{ctp-red}{完整的官学体系}，也是第一个官学体系。中央专设“六学”：国子学、太学、四门学、书学、算学、律学，
前三学为\textcolor{ctp-red}{大学性质}后三学为\textcolor{ctp-red}{专科学校性质}。地方官学分为\textcolor{ctp-red}{经学、医学、崇玄学}三级。

\paragraph{私学}
私学是由\textcolor{ctp-red}{私人举办}的学校，始于春秋时期，打破\textcolor{ctp-red}{“学在官府”}的局面。

\textcolor{ctp-red}{孔子}所办私学\textcolor{ctp-red}{时间长久，规模宏大，成绩突出，对后世影响深远}。首开平民教育，主张“有教无类”，开创了私学的先河。
孔子认为教育的主要目的是培养人良好的道德品质，并将教育内容划分为四科：\textcolor{ctp-red}{德行，言语，政事，文学}。

\textcolor{ctp-red}{初级私学称为蒙学}，高级私学是对具有一定文化基础的成年人进行专门教育的学校。私学\textcolor{ctp-red}{没有固定的年限}，
每个私学通常只有一个老师，采用\textcolor{ctp-red}{个别教学方法}。

\paragraph{书院}
唐宋以后，书院成为重要的教学机构。\textcolor{ctp-red}{白鹿洞书院}是中国古代最著名的书院之一。
此书院的学规\textcolor{ctp-red}{《白鹿洞书院揭示》}由南宋理学家\textcolor{ctp-red}{朱熹}编订，成为中国古代书院学规的重要范本，对后世书院的发展产生了深远影响。
宋朝书院已经出现官学化倾向，政府加强对书院的控制，书院逐渐纳入官学体系。

明朝\textcolor{ctp-red}{东林书院}“家事国事天下事，事事关心”，书院不仅是一个重要的\textcolor{ctp-red}{文化学术中心}，
也是一个重要的\textcolor{ctp-red}{政治活动中心}。书院具有如下主要特点：
\begin{itemize}
  \item 教育与科研密切结合
  \item 重视学术争辩和切磋，开展不同学派间的交流
  \item 重视学术自学，提倡独立研究
  \item 课程设置简约，学生可根据兴趣爱好学习
  \item 师生关系融洽
\end{itemize}

\paragraph{学校教育内容} 六艺：礼、乐、射、御、书、数。
\paragraph{选士制度} 两汉时期察举制，魏晋南北朝时期九品中正制，隋唐科举制。

\subsubsection{中国近代教育}
\paragraph{新式学堂} 1862年创办的\textcolor{ctp-red}{京师同文馆}是中国近代重要的新式学堂之一，主要培养翻译人才。
\paragraph{近代学制}
\subparagraph{癸卯学制}又称“奏定学堂章程”，“1904年学制”，是\textcolor{ctp-red}{中国近代第一部颁布并实施的学制}。
其诞生标志着中国近代教育走向\textcolor{ctp-red}{制度化、法制化}阶段。
\begin{cquote}
  癸卯学制规定了小学、初级中学、高级中学、大学四个层次的学制体系，明确了各级学校的学制年限和课程设置，为中国近代教育的发展奠定了基础。
\end{cquote}

\subparagraph{壬子癸丑学制}又称“1912年学制”，辛亥革命后，在\textcolor{ctp-red}{蔡元培}的主持下制定的\textcolor{ctp-red}{民国新学制}，
是中国近代第一个\textcolor{ctp-red}{资产阶级}性质的学制。 确定了女子受教育的权利。

\subparagraph{壬戌学制}又称“1922年学制”，“新学制”，“六三三学制”是\textcolor{ctp-red}{中国近代第二个资产阶级学制}，
由\textcolor{ctp-red}{蔡元培}主持制定。是中国近代实施时间最长、影响最大、也是最成熟的学制。壬戌学制的主要特点有：
\begin{enumerate}
  \item 缩短小学毕业年限，延长中学修业年限
  \item 若干措施注意根据地方实际需要，不做硬性规定
  \item 重视学生的职业教育和补习教育
  \item 课程和教材内容侧重实用
  \item 实行选科制和分科教育，兼顾学生升学和就业两种准备
\end{enumerate}

\paragraph{近代教育行政机构}
\begin{itemize}
  \item 1905年12月设立学部，最高长官为尚书，首任为荣庆
  \item 辛亥革命后改名为教育部，最高长官称教育总长，首任蔡元培
  \item 1906年4月，设立省级教育机构
\end{itemize}

\subsubsection{建国后的教育}
\paragraph{改造旧教育，创建新教育} 1949年11月1日，中央人民政府教育部成立，马叙伦任部长。进行了以下举措：
\begin{enumerate}
  \item 改造旧学校。政府逐步接管私立学校，改造为“国家办学”的单一教育体制。
  \item 颁布新学制，1951年10月确定新学制，小学五年一贯制、中学三三制、大学和专门学院4-5年。“学校向工农“，“教育为国家建设服务”。
  \item 学习苏联教育经验，1952年起开展学习苏联教育热潮，影响最大著作为凯洛夫的《教育学》。
\end{enumerate}

\paragraph{新时期的教育改革与发展}
“教育要面向现代化，面向世界，面向未来”——邓小平 \\
1985年后，\textcolor{ctp-red}{“三个面向”}被载入党和国家的重要决定和文件中，成为长期指导我国教育改革发展的战略指导方针。

\noindent
1977年5月，\textcolor{ctp-red}{“科教兴国”}战略提出，强调科技和教育是国家发展的重要支柱。 \\
党的十三大“百年大计，教育为本”提出，强调教育在国家发展中的基础地位。 \\
对马克思主义教育理论的继承和发展。

1993年2月13日，\textcolor{ctp-red}{“教育改革和发展纲要”}发布，提出了教育改革的总体目标和主要任务。
\textcolor{ctp-red}{素质教育称为新时期教育改革的关键词和风向标}。

1991年1月《面向21世纪教育振兴行动计划》\textcolor{ctp-red}{“跨世纪素质教育工程”}提出，强调全面提高国民素质，培养适应21世纪需要的人才。

2010年7月29日，《国家中长期教育改革和发展规划纲要（2010-2020年）》发布，提出了未来十年教育改革和发展的总体目标和主要任务。

2019年2月23日，\textcolor{ctp-red}{《中国教育现代化2035》}发布，是我国第一个以\textcolor{ctp-red}{教育现代化为主题的中长期战略规划}，
是新时代推进教育现代化、建设教育强国的纲领性文件。

\subsection{第三节 \space 西方教育的发展}

\subsubsection{西方古代教育}
\paragraph{古希腊教育} 古希腊教育是西方教育的开端。\\
斯巴达“尚武”教育，强调军事训练和纪律，教育全部由国家组织管理和控制；雅典“尚文”教育，注重文化、艺术和哲学的培养，既有公共教育，也有私人教育。

\paragraph{古罗马教育} 主要分为\textcolor{ctp-red}{共和时期教育、帝国时期教育}两大阶段。

共和国时期，罗马经济以农业为主，此时教育为农民——军人教育为主要特征。具有希腊语学校和拉丁语学校两种学校系统。

\paragraph{西欧中世纪教育} “僧侣获得知识教育的垄断地位”，教会学校是几乎唯一的教育机构。

僧侣是主要的教育者，基督教教育是西欧教育的主干。学校分为修道院学校、主教学校和教区学校。修道院学校的教师由教士担任，教学方法以口授为主。
一所主教学校只有一位教师，通常由主教担任。

西欧中世纪早期，世俗封建主教育的形式主要是\textcolor{ctp-red}{骑士教育}，这是一种特殊形式的\textcolor{ctp-red}{家庭教育}。

西欧中世纪晚期，\textcolor{ctp-red}{中世纪大学和城市学校}出现。最早的大学是\textcolor{ctp-red}{意大利萨莱诺大学}，
随后是\textcolor{ctp-red}{巴黎大学}和\textcolor{ctp-red}{牛津大学}等。

城市大学包括工业行会开办的行会学校，如商人联合会设立的\textcolor{ctp-red}{基尔特学校}。

\paragraph{西方近代教育}
\subparagraph{近代德国教育}
\textcolor{ctp-red}{赫尔巴特}的\textcolor{ctp-green}{《普通教育学》}是世界上第一部具有科学体系的教育学著作。
德国学前教育学家\textcolor{ctp-red}{福禄贝尔}创办的\textcolor{ctp-red}{幼儿园}影响世界学前教育发展。

\vspace{0.6em}

\textcolor{ctp-lavender}{近代西方的教育大多起源于德国，并对其他国家产生了重要影响。}

\subparagraph{19世纪以后西方教育的主要特点}
\begin{enumerate}
  \item 国家重视教育，建立公共教育体系
  \item 普及初等义务教育
  \item 实科教育迅速发展
  \item 重视教育立法，依法治教
\end{enumerate}

\subsection{第四节 \space 中国未来教育的发展趋势}
\begingroup
\setstretch{1.05}
\setlength{\columnsep}{1em}
\columnratio{0.50,0.50}
\begin{paracol}{2}

  \begin{enumerate}
    \item 信息技术在教育中应用更加广泛
    \item 教育培养目标转向能力培养为主
    \item 混合式学习更加普遍
    \item 学生的培养将更加个性化
    \item 学习更以学生为中心
  \end{enumerate}

  \switchcolumn

  \begin{enumerate}
    \item [6] 教师的角色和作用将发生变化
    \item [7] 学校的办学模式将发生改变
    \item [8] 终身学习将成为人们的生活方式
    \item [9] 教育对外交流与合作将更进一步发展
    \item [10] 民办教育将继续发展
  \end{enumerate}

\end{paracol}
\endgroup

\section{第二章 教育思想的演进}
\subsection{第一节 中国教育思想的演进}

\subsubsection{中国古代教育思想}
\paragraph{孔子的教育思想} 孔子是中国古代最伟大的教育家和思想家之一，他的教育思想对中国文化和教育产生了深远的影响。

\textcolor{ctp-red}{孔子的教育思想集中体现在《论语》中，主要包括“性相近，习相远”、“有教无类”、“因材施教”等核心观点。}

“学而优则仕”是孔子教育思想的核心内容之一，强调通过学习和修养来提升个人的品德和能力，从而为社会做出贡献。

孔子最早提出\textcolor{ctp-peach}{启发式教学}，“不愤不启，不悱不发”，强调通过启发和引导学生自主思考和探索知识。
\textcolor{ctp-red}{当学生经过认真思考并达到一定程度时，教师要恰到好处地进行启发和开导。}孔子同时也是最早采用因材施教方法的教育家。

\paragraph{墨子的教育思想} 与孔子创立的儒家学派并称为\textcolor{ctp-red}{“显学”}。有非儒即墨之说。

\textcolor{ctp-red}{素丝说}，“染于苍则苍，染于黄则黄”，强调环境对人的影响，主张通过教育来改变人的性格和行为。
\textcolor{ctp-red}{兼爱非攻}，主张无差别的爱和反对战争，强调教育的社会责任和道德价值。

\paragraph{《学记》的教育思想}
\textcolor{ctp-red}{《学记》}是中国古代也是世界上最早的一篇专门论述教育问题的论著，被称为\textcolor{ctp-red}{教育学的雏形}。
是先秦儒家教育经验的初步总结，\textcolor{ctp-red}{是中国古代教育思想开始从其他学科独立出来的一个重要标志}。

《学记》的\textcolor{ctp-red}{四大原则}：\textcolor{ctp-red}{预防性原则、及时施教原则、循序渐进原则、学习观摩原则}。
实施教学不超越次序叫做循序渐进，同学间相互学习帮助叫做观摩切磋。

《学记》\textcolor{ctp-red}{“善喻”}教学法：
\begin{itemize}
  \item \textcolor{ctp-red}{“道而弗牵”} 教师起主导作用，引导学生主动地投入学习，而不能牵制束缚学生的思维。
  \item \textcolor{ctp-red}{“强而弗抑”} 教师在学生遇到困难时给予适当的帮助，但不压制学生的自主性。
  \item \textcolor{ctp-red}{“开而弗达”} 教师应启发学生思考，引导学生自主探索知识，而不是直接给出答案。
\end{itemize}

\textcolor{ctp-red}{“学者有四失”}：“或失则多，或失则寡，或失则易，或失则止”

\paragraph{《大学》的教育思想} 是一篇\textcolor{ctp-red}{道德教育的专论。}
\textcolor{ctp-peach}{“皆以修身为本”}

\noindent
\textcolor{ctp-red}{“三纲领”}：明明德、亲民、止于至善。\\
\textcolor{ctp-red}{“八条目”}：格物、致知、诚意、正心、修身、齐家、治国、平天下。

\paragraph{《颜氏家训》的教育思想} 颜之推所著，“古今家训之祖”。
\begin{enumerate}
  \item 重视儿童的早期教育，认为一个人的发展，幼儿时期是奠定基础的重要阶段。
  \item 在家庭中对儿童进行教育，应当严慈相济，反对溺爱和放任。
  \item 多子女家庭，对子女应一律平等。
  \item 语言是儿童学习的重要内容，对儿童进行的语言教育应注意规范。
  \item 必须重视道德教育，对儿童进行道德教育，不应该是说教式的。
\end{enumerate}

\textcolor{ctp-red}{提出勤学，切磋，眼学}三大学习主张。

\paragraph{韩愈的教育思想} 韩愈是唐代著名文学家，思想家和教育家。教育思想主要体现在《师说》《进学解》中。

教师的职责主要有三项：传道受业解惑。教师的职责是传授知识，指导学生学习，解答学生的疑问。

“闻道有先后，术业有专攻”，强调不同学科有不同的专业领域，学生应该根据自己的兴趣和特长选择适合自己的学科进行学习。

“弟子不必不如师，师不必贤于弟子”，强调教师和学生之间的关系是平等的，学生也可以在某些方面超过教师。

在《进学解》中论及了\textcolor{ctp-red}{“教学艺术”}，主要观点有：勤学深思，博学求精，把学习与独创结合起来，注重实践和应用。

\paragraph{朱熹的教育思想} 朱熹是南宋著名的理学家和教育家，他的教育思想主要体现在《朱子语类》中。

最值得关注的是\textcolor{ctp-red}{论述“小学”和“大学”教育}和“朱子读书法”，读书法有六条：
\begin{enumerate}
  \item 循序渐进：按照一定的顺序进行学习，先学基础知识，再逐步深入。
  \item 熟读精思：通过反复阅读和深入思考来理解和掌握知识。
  \item 虚心涵泳：保持谦虚的态度，不要先入为主，牵强附会，反复咀嚼，细心玩味。
  \item 切记体察：必须见之于自己的实际行动，身体力行。
  \item 着紧用力：读书必须抓紧时间，发愤忘食，必须精神抖擞，勇猛奋发。
  \item 居敬持志：保持敬畏之心和坚定的志向，专注于学习和追求真理。
\end{enumerate}

\noindent
\textcolor{ctp-red}{博学、审问、慎思、明辨、笃行}是教学和学习的五个阶段。

\subsubsection{中国近代教育思想}
\paragraph{蔡元培的教育思想} 五育并举：
\textcolor{ctp-rosewater}{
  \begin{enumerate}
    \item 军国民教育主张将军事教育引入学校社会教育。
    \item 实利注意教育，以人民生计为普通教育中坚。
    \item 公民道德教育是以西方资产阶级的自由、平等、博爱思想为教育内容。
    \item 世界观教育。
    \item 美感教育，与世界观教育紧密相连。
  \end{enumerate}
}

\noindent
蔡元培还提出：\textcolor{ctp-red}{以美育代替宗教}以及\textcolor{ctp-red}{思想自由，兼容并包}的教育方针。

\paragraph{陶行知的教育思想} 三个最基本的教育命题是：
\textcolor{ctp-red}{“生活即教育”“社会即学校”“教学做合一”}

\noindent
\textcolor{ctp-red}{陶行知是中国创造教育的先驱}，被誉为“伟大的人民教育家”“万世师表”。

\paragraph{陈鹤琴的教育思想} 核心是\textcolor{ctp-red}{活教育}
\begin{enumerate}
  \item \textcolor{ctp-red}{目的论} 做人，做中国人，做现代中国人
  \item \textcolor{ctp-red}{课程论} 符合儿童活动和生活方式，符合儿童与自然、社会的交往方式
  \item \textcolor{ctp-red}{方法论} 做中教、做中学、做中求进步
\end{enumerate}

\paragraph{黄炎培的教育思想} \textcolor{ctp-red}{职业教育之父}

\subparagraph{大职业教育} \textcolor{ctp-maroon}{办学宗旨社会化、培养目标社会化、学制社会化、办学方式社会化、办学过程社会化}

职业道德的基本要求概括为\textcolor{ctp-red}{敬业乐群}

\paragraph{晏阳初的教育思想} \textcolor{ctp-red}{国际平民教育之父}

推行教学\textcolor{ctp-red}{三大方式}：学校式、家庭式、社会式。

\subsection{第二节 西方教育思想的演进}
\subsubsection{西方古代教育思想}
\paragraph{苏格拉底的教育思想} 苏格拉底是古希腊著名的哲学家和教育家，他的教育思想主要体现在他的教学方法和教育理念中。
\textcolor{ctp-red}{苏格拉底的教育方法主要是通过对话和提问来引导学生思考和探索知识，这种方法被称为“苏格拉底式教学法”。}

\paragraph{柏拉图的教育思想} 柏拉图是古希腊著名的哲学家和教育家，他的教育思想主要体现在他的著作\textcolor{ctp-red}{《理想国》}中。
\paragraph{亚里士多德的教育思想} 第一位划分年龄分期，并提出教育主张的教育家。

\subsubsection{西方近现代教育思想}
\paragraph{夸美纽斯的教育思想} “教育科学真正的奠基人”

\textcolor{ctp-red}{《母育学校》} 西方教育史上第一部详细记录幼儿教育的学前教育专著。

\textcolor{ctp-red}{《大教学论》} 是西方教育史上第一部成体系的教育学著作。

提出许多适应于自然的原则，如直观性原则、循序渐进和系统性原则、自觉性和主动性原则、量力性原则、因材施教原则。

\paragraph{卢梭的教育思想} 代表作\textcolor{ctp-red}{《爱弥儿》}，是西方教育史上第一部系统论述自然主义教育的著作。

\textcolor{ctp-red}{自然教育理论}，自然教育的目的是培养自然天性充分得到发展的\textcolor{ctp-red}{自然人}。

\paragraph{赫尔巴特的教育思想} 代表作\textcolor{ctp-red}{《普通教育学》}，是世界上第一部学科形态的教育学，其出版标志着规范教育学的建立。
“形式阶段”理论\textcolor{ctp-red}{明了，联想，系统，方法}四个阶段。

\paragraph{福禄贝尔的教育思想} \textcolor{ctp-red}{近现代学前教育理论的奠基人}，“幼儿园之父”。创办了世界上第一所幼儿园。
\textcolor{ctp-red}{恩物游戏}，通过游戏来促进儿童的全面发展，游戏是人在早期发展阶段最重要的精神产物。

\paragraph{杜威的教育思想} \textcolor{ctp-red}{教育即生活、学校即社会、从做中学}

\textcolor{ctp-red}{现代三中心}
\begin{enumerate}
  \item 儿童中心：主张在学校生活中，“儿童是起点，是中心，而且是目的”，强调以儿童为中心的教育理念。
  \item 经验中心：要以经验作为学生的学习内容，从经验中学，不断丰富他们的经验，增强不断变化世界的适应能力。
  \item 活动中心：教学过程应该以儿童的实践活动为中心，“从做中学”，教学过程就是“做”的过程。
\end{enumerate}

\section{课程}
\subsection{第一节 课程概述}
\paragraph{课程的具体表现形式} \textcolor{ctp-red}{课程计划、课程标准、教材}

\subsection{第二节 课程理论}
\paragraph{学科中心课程论} 学校课程应以学科的分类为基础，\textcolor{ctp-red}{以学科教学为中心}，
以掌握学科的基本知识、基本规律和相应的技能为目标。

\paragraph{学科中心课程论的发展} 早期发展人物是\textcolor{ctp-red}{斯宾塞}——《什么知识最有价值》（1859）

\textcolor{ctp-red}{赫尔巴特}对学科中心课程论做出重要贡献，他在\textcolor{ctp-red}{普通教育学}提出，
根据经验的兴趣和思辨的兴趣分别设置课程。

\textcolor{ctp-red}{布鲁纳}使当代学科中心课程论进一步发展。

\textcolor{ctp-red}{要素主义课程}，\textcolor{ctp-peach}{课程要以共同的、不变的文化要素为核心}。
\textcolor{ctp-red}{永恒主义课程}也是学科中心论的主要流派，提出课程应该是\textcolor{ctp-red}{不变的学问}（读，写，算）

\paragraph{学科中心论的优缺点} 需要辩证看待学科中心论的优缺点，不能片面强调其优点而忽视其缺点，也不能片面批评其缺点而否定其优点。
\begingroup
\setstretch{1.05}
\setlength{\columnsep}{1em}
\columnratio{0.50,0.50}
\begin{paracol}{2}
  \textcolor{ctp-red}{学科中心论的优点}
  \begin{enumerate}
    \item 按学科组织教学内容，有利于文化遗产的保存与传递
    \item 按学科教授相关知识，有利于系统知识的掌握
    \item 根据学科逻辑编排教材，有利于促进学生思维的发展
    \item 课程主要任务清晰，有利于基础知识与基本技能的传授
    \item 课程的构成比较单纯，容易教学和评价
  \end{enumerate}

  \switchcolumn

  \textcolor{ctp-red}{学科中心论的缺点}
  \begin{enumerate}
    \item 按学科组织教学内容，容易把相关知识割裂开来
    \item 学科内容往往与社会关心的问题及发生的事件相脱节
    \item 对学生的兴趣和需要重视不够
  \end{enumerate}
\end{paracol}
\endgroup

\subsubsection{活动中心课程论}
\paragraph{活动中心课程的基本观点} \textcolor{ctp-red}{又称儿童中心课程论}，其主要代表有\textcolor{ctp-red}{杜威、卢梭}等。

\paragraph{活动中心课程论的发展} 法国思想家卢梭，\textcolor{ctp-red}{“发现儿童”}是其在教育思想史上最伟大的贡献。
\textcolor{ctp-red}{杜威}是活动中心课程论的主要代表人物，他在《民主主义与教育》中提出了活动中心课程的基本观点。

\paragraph{活动中心课程论的优缺点}

\textcolor{ctp-red}{活动课程具有如下优点：}
\begin{enumerate}
  \item 重视学生的需要与兴趣，尊重学生的主体性，有利于学生学习的主动性、积极性的发挥。
  \item 强调教材的心理组织，有利于学生在与文化，与科学知识的交互作用的过程中，获得人格的不断发展。
  \item 强调实战活动，重视学生通过亲身体验获得直接经验，有利于培养学生解决实际问题的能力。
  \item 重视课程的综合性，主张以社会生活问题来总和各种知识，有利于学生获得对世界的完整认识。
\end{enumerate}

\textcolor{ctp-red}{活动课程的缺点：}
\begin{enumerate}
  \item 过分地夸大了儿童个人经验的重要性，忽视系统的学科知识的学习，忽视儿童思维力和其他智力品质的发展。
  \item 相对学科课程而言，活动课程难以设计与组织实施。
  \item 降低了教师的指导作用，教师是儿童学习的参谋或顾问，帮助学生解决疑难问题，容易使教师丧失责任心。
\end{enumerate}

\subsubsection{社会中心课程论}
\paragraph{社会中心课程论的基本观点} 学科中心课程以学科为中心，活动中心课程以儿童为中心，而\textcolor{ctp-red}{社会中心课程理论}则以社会问题为中心。

\subsection{第三节 课程的分类}
\subsubsection{按照学科内容的属性分类}
按照学科内容的属性来划分，可以把课程分为\textcolor{ctp-red}{学科课程和经验课程}。

\paragraph{学科课程}具有\textcolor{ctp-lavender}{科目性、结构性、专业性、系统性、预设性}等特点。

学科课程的主要优点是：有助于人类文化遗产的完整保存与传递；有助于学习获得者获得系统连贯的文化科学知识；有助于教学的组织、评价以及教学效率的提高。

学科课程的主要缺点有：容易导致儿童机械被动学习，或使儿童仅仅称为学习活动的旁观者；容易忽略必要的个人知识、实践知识的学习和获取；
过分强调学科课程，容易导致儿童完整的生活世界被割裂和肢解；过分强调学科课程的学习，容易导致学校与现实社会生活的隔离。

\paragraph{活动课程} “进步教育运动”的代表人物是杜威。

活动课程的主要优点有：课程贴近儿童当前生活；有助于儿童学习和掌握必要的个人实践知识，促进能力发展；儿童是课程活动积极主动的参与者和建构者；
学习过程是儿童全人格整体参与的过程，是理智与情感的融合过程，是思维与行动的统一过程。

活动课程的主要缺点有：不能给学习者提供系统的文化科学知识；在实施过程中，儿童思维能力的发展容易落空；活动课程的组织要求较高。

\subsubsection{根据课程内容的组织方式分类}
按照课程内容的组织方式来划分，可以把课程分为\textcolor{ctp-red}{分科课程、综合课程}。

分科课程是\textcolor{ctp-red}{从不同门类的学科中选取知识，按照知识的逻辑体系，以分科教学的形式向学生传授知识的课程。}

综合课程又分为：\textcolor{ctp-red}{简单相加式综合、相关课程综合、一体化课程}。

\subsubsection{根据课程的管理归属分类}
按照课程的管理归属来划分，可以把课程分为\textcolor{ctp-red}{国家课程、地方课程、校本课程}。

国家课程具有\textcolor{ctp-red}{强制性，统一性、公共性、基础性}等特点。

校本课程具有\textcolor{ctp-red}{校本性、多样性、灵活性、个性化、特色化、自主性、动态性}等特点。

\subsubsection{根据课程的修习要求分类}
按照课程的修习要求来划分，可以把课程分为\textcolor{ctp-red}{必修课程、选修课程}。

\subsubsection{根据课程呈现形式划分}
按照课程的呈现形式来划分，可以把课程分为\textcolor{ctp-red}{显性课程和隐性课程}。

隐性课程具有如下特点：\textcolor{ctp-red}{普遍性、隐蔽性、非计划性、渗透性、感染性、长效性、动态性}等特点。

\section{第四章 教学}
\subsection{第一节 \space 教学概述}
\subsubsection{教学的概念}
\textcolor{ctp-red}{教学是学校教育工作的一个重要组成部分，是学校教育的主要途径。}

\subsubsection{教学的任务}
\textcolor{ctp-rosewater}{
  \begin{enumerate}
    \item 引导学生掌握科学文化基础知识和基本技能
    \item 发展学生的体力、智力、创造力和实践精神
    \item 培养学生的道德品质和审美情趣
    \item 发展学生的个性
  \end{enumerate}
}

\subsubsection{教学的意义}
\textcolor{ctp-red}{教学是实现教育目的的基本途径}

\textcolor{ctp-red}{教学是学校工作的中心环节}
\textcolor{ctp-peach}{教学是学校的中心工作，学校工作必须坚持以教学为中心。}

\subsection{第二节 \space 教学过程}
\subsubsection{教学过程的实质}
\paragraph{教学是一种特殊的认识过程} 具有\textcolor{ctp-red}{间接性、引导性、简捷性}。

\subsubsection{教学过程的基本规律}
\noindent
\textcolor{ctp-red}{
  获取直接经验与学习间接经验相统一 \\
  掌握知识与发展能力相统一 \\
  传授知识与思想品德教育相统一 \\
  教师主导作用与学生主体作用相统一
}

\subsection{第三节 \space 教学环节与组织形式}
\subsubsection{教学环节}
\textcolor{ctp-sky}{备课、上课、作业的布置与批改、课外辅导、学生学业成绩的检查、评定}

\paragraph{备课} 根据课程标准的要求和本门课的特点，结合学生的具体情况，选择最合适的表达方法和顺序。
需要做好\textcolor{ctp-red}{备教材、备学生、备教法}三方面的准备工作。

\paragraph{上课} 上课是整个教学工作中的中心环节，是提高教学质量、培养学生的关键。
\textcolor{ctp-red}{教学目标要确定、教学内容要正确、教材组织要恰当、教学方法要得当、教学结构要紧凑、教学效果要良好。}

\paragraph{作业的布置与批改} 作业布置与批改是教学工作的重要组成部分。

\begin{enumerate}
  \item 作业内容要符合学科课程标准和教科书的要求
  \item 作业分量要适当、难易要适度、作业时间要有所控制
  \item 作业目的性要明确，要求要具体
  \item 作业应与教科书的内容有逻辑关系，但不是照搬
  \item 作业指导要讲方法和策略
  \item 要及时检查和批改作业
\end{enumerate}

\subsubsection{教学组织形式}
\noindent
\textcolor{ctp-red}{班级授课制是目前学校最普遍的教学组织形式。} \\
\textcolor{ctp-red}{分组教学容易出现快组学生容易骄傲，普通组、慢组学生的学习积极性普遍降低。} \\
个别教学是适应学生个别差异，发展学生个性的教学。

\subsection{第四节 \space 教学方法}
\subsubsection{常用教学方法}
有\textcolor{ctp-red}{讲授法、谈话法、讨论法、实验法、实习作业法}、练习法、读书指导法、演示法、\textcolor{ctp-red}{研究法}。

\textcolor{ctp-red}{讲授法有讲述、讲解、讲读、讲演}四种方式。

谈话法又称为\textcolor{ctp-red}{问答法}，是在学生已有知识经验的基础上提出问题，引导学生积极思考。
可分为\textcolor{ctp-red}{启发性谈话、问答式谈话、指导总结谈话}三种。

讨论法是在教师的指导下，学生围绕某一问题进行交流和探讨的一种教学方法。讨论法\textcolor{ctp-red}{以学生活动为主，以教师指导为辅。}

实验法和研究法的区别在于，\textcolor{ctp-red}{实验法侧重于获取知识、巩固知识和培养操作技能，研究法侧重于创造性地解决问题并独立探索知识。}

实习作业法是将书本知识运用于实践的一种教学方法。

\section{第五章 \space 课堂管理}
\subsection{第一节 \space 课堂管理概述}
\begin{cquote}
  秧田式座位排列的优点是教室整齐划一、空间利用充分，有利于教师管理课堂、维持秩序和有计划地传授知识，
  有利于教师观察学生的一切活动。
\end{cquote}

\subsubsection{课堂气氛的营造}
\begin{table}[!ht]
    \centering
  \small
  \setlength{\tabcolsep}{3pt}
  \renewcommand{\arraystretch}{1.15}
  \begin{tabular}{@{}p{0.12\textwidth}p{0.25\textwidth}p{0.25\textwidth}p{0.30\textwidth}@{}}
    \hline
        ~ & 积极的 & 消极的 & 对抗的 \\ \hline
        注意状态 & 注意稳定集中，全神贯注，如痴如醉 & 学生呆若木鸡，老师无精打采 & 学生注意力不在学习，老师常被迫维护课堂或放任自流 \\ \hline
        情感状态 & 积极愉快，情绪饱满 & 学生压抑无精打采，老师缺少情感和责任心 & 学生有意捣乱敌视教师，教师冷漠或大发脾气 \\ \hline
        意志状态 & 信心十足，坚韧不拔 & 学生害怕困难，老师不愿引导，知难而退 & 学生缺乏信心设法逃避，老师对学生失望或厌弃学生 \\ \hline
        定式状态 & 学生确信教师，教师确信为学生服务 & 学生对老师课的正确性、方法等无所谓，老师漫不经心 & 学生怀疑不信任教师，教师漠不关心或厌烦学生 \\ \hline
        思维状态 & 学生思维活跃，积极学习，教师富有灵感。 & 学生被动思维，不愿动脑，老师不假思索、敷衍了事 & 学生不思问题，不动脑筋，老师抱怨学生愚笨 \\ \hline
    \end{tabular}
\end{table}

\textcolor{ctp-red}{责罚学生时需避免涟漪效应}

\paragraph{良好课堂气氛的营造策略} 课堂气氛是课堂中各种因素相互作用的结果。

\textcolor{ctp-red}{正确分析与鉴定课堂气氛； 以积极的情感感染学生； 抓典型、树榜样、立威信； 妥善处理矛盾冲突。}

\subsection{第三节 \space 课堂管理的原则}
\noindent
\textcolor{ctp-red}{尊重学生；民主性；时效性；方向性；效率性；适度性；主体性}

\subsection{第四节 \space 课堂问题行为处理}
\subsubsection{课堂问题行为及其类型} 
具有四个特征： \textcolor{ctp-red}{普遍性、以轻度为主、消极行为、可以矫正}。

\subsubsection{课堂问题行为矫正方法}
常用矫正方法包括\textcolor{ctp-red}{行为矫正和认知矫正}两大类。

行为矫正有\textcolor{ctp-red}{奖励、惩罚、负强化}等方法。

\subsubsection{课堂问题行为的矫正程序} 
具有六个基本环节：\textcolor{ctp-red}{觉察，诊断，目标，改正，检评，追踪}。

\section{第六章 \space 教育评价}
\subsection{第一节 \space 教育评价概述}
\subsubsection{教育评价的概念}
\textcolor{ctp-blue}{教育评价是对教育活动的价值进行判断和评价的过程。}
教育史上著名的\textcolor{ctp-red}{“八年实验研究”}，其成果体现在《史密斯-泰勒报告》中。

\subsubsection{教育评价的功能}
\textcolor{ctp-red}{教育评价的功能主要有：教育、导向、改进、激励、鉴定、管理}等功能。

\subsection{第二节 \space 教育评价的类型}
\paragraph{量化评价与质性评价} 质性评价通常表现为书面的“鉴定”或“评语”。

\subsection{第三节 \space 教育评价的原则}
\begin{enumerate}
  \item 效用评价
  \item 定性评价与定量评价相结合
  \item 静态评价与动态评价相结合
  \item 他人评价与自我评价相结合
\end{enumerate}

\end{document}